\section{Background and Related Work}
\label{sec:related}

\subsection{Intelligent Textbooks and Learning Graphs}

Intelligent textbooks have a history reaching back to adaptive hypermedia in the early 1990s. Brusilovsky, Sosnovsky, and Thaker trace several generations, from hand-built domain models to interaction data mining, and note that the line of work never stopped even as attention to it waxed and waned \cite{brusilovsky2022return}. The book evaluated here uses a simple five-level scale: static content, interactive content (hyperlinks, quizzes, simulations, glossary), adaptive content, chatbot integration, and autonomous AI that generates lessons for each reader. The books the benchmark generates are Level~2.

The structural device that makes generation tractable is the \textit{learning graph}: a directed acyclic graph of concepts in which an edge means ``must be understood before.'' The idea sits in a long line of work on prerequisite structure, from knowledge spaces, in which a prerequisite partial order defines the states of knowledge a learner can be in \cite{doignon1999knowledge}, to concept maps, in which concepts joined by labeled links form propositions around a focus question \cite{novak2008concept}. A learning graph also supports a mechanical quality check, since a valid graph has no cycle, and a mechanical ordering, since any topological order is a legal chapter order. Learning outcomes and quiz items are labeled with the revised Bloom taxonomy \cite{anderson2001taxonomy}, and the course description that seeds the pipeline follows backward design: results first, then evidence, then lessons \cite{wiggins2005ubd}.

Interactive elements in the pipeline are MicroSims: small, browser-based simulations, each tied to one concept, embedded by iframe, and specified in a standard layout so that a language model can generate and revise them \cite{mccreary2025microsims}. This paper does not evaluate MicroSim generation; it is the stage the benchmark deliberately leaves out (Section~\ref{sec:design}).

\subsection{Small Language Models and Local Inference}

``Small'' has no fixed meaning. Surveys of small language models typically concern models of 100 million to a few billion parameters \cite{lu2024slm, wang2024slmsurvey}. Belcak et~al.\ instead define an SLM functionally, as a model that can run on a common consumer device with latency low enough to serve one user's agentic requests \cite{belcak2025slm}. We follow the functional reading, as does the book: an SLM is \textit{a model small enough to run on one machine you control}. On the machine tested here, that admits models of 24 to 128 billion total parameters, many of which are sparse mixtures of experts \cite{shazeer2017moe} that activate only 3 to 5 billion parameters per token; \texttt{gpt-oss-120b}, for example, has about 117 billion parameters of which 5.1 billion are active for each token \cite{openai2025gptoss}. This is large by the standard of the surveys and small beside frontier systems, and we use ``SLM'' in this operational sense throughout.

Three technical facts make the 128~GB class workable. Weight quantization shrinks a model to a fraction of its 16-bit size with modest quality loss \cite{dettmers2022llmint8, frantar2022gptq}; every model tested here ran at its publisher's default 4-bit file. Sparse mixtures of experts read only the active experts for each token, which raises decode speed on bandwidth-limited hardware. And unified memory lets weights and conversation context share one pool, so that a long context competes directly with the weights rather than spilling to a separate device memory \cite{nvidia2025dgxspark}. Inference servers differ in what they optimize: llama.cpp and the Ollama server built on it favor simple local deployment \cite{llamacpp, ollama}, while vLLM favors throughput under concurrent load through paged attention \cite{kwon2023vllm}. This evaluation used Ollama and left vLLM as a second-pass candidate.

\subsection{Agent Harnesses, Skills, and Tool Calling}

An agent harness runs a loop in which the model proposes an action, the harness executes it through a tool, and the result is appended to the conversation \cite{yao2023react, anthropic2024agents}. Tool calling is a learned skill of the model and has its own benchmarks \cite{patil2023gorilla, patil2025bfcl}; agentic benchmarks show that poor long-horizon reasoning and instruction following are the main obstacles for open models acting as agents \cite{liu2023agentbench}. Software-engineering benchmarks test agents on real repositories \cite{jimenez2024swebench}. Agent Skills package instructions, scripts, and reference files in a folder with a single \texttt{SKILL.md} entry point that an agent loads on demand \cite{anthropic2025skills}; the pipeline benchmarked here is a chain of such skills. The harness is Claude Code \cite{claudecode}. That choice is consequential, because a harness written for one model family can disadvantage another; we return to this confound in Section~\ref{sec:discussion}.

\subsection{Evaluating Language Models and Generated Text}

Evaluation practice supplies the vocabulary this paper borrows. HELM argues for standardized conditions, multiple metrics, and released raw outputs \cite{liang2022helm}. The lm-evaluation-harness experience shows how sensitive scores are to setup details and how often comparisons are not reproducible \cite{biderman2024lessons}. Reproducibility checklists \cite{pineau2021reproducibility}, work on reporting variance across runs \cite{dodge2019show}, and the finding that non-determinism can hide whether an improvement is real \cite{henderson2018rl} motivate our insistence that a single run supports a statement about that run only. Using a language model to grade generated text is attractive and carries documented biases: position, verbosity, and self-enhancement \cite{zheng2023judge}, and a preference for one's own generations \cite{panickssery2024selfpref}. Long prompts are a known weak point of language models \cite{liu2024lost, hsieh2024ruler}, which matters here because a harness resends a growing conversation on every turn. Model cards recommend stating intended use and limits \cite{mitchell2019modelcards}, which is the posture of our Section~\ref{sec:discussion}. Finally, performance-evaluation practice has long catalogued common mistakes, such as unrepresentative workloads and biased goals \cite{jain1991art}, and the A/A test of online experimentation, in which identical variants validate the apparatus \cite{kohavi2020trustworthy}, motivates our control runs.

\subsection{Gap}

Prior work evaluates models on short tasks, on tool calling in isolation, or on software repositories. We did not find a published evaluation of open-weight models on a single desktop-class machine performing an end-to-end, multi-stage, document-generation workflow under a skill-following harness, scored without a model as judge, and accompanied by a worked method for repeating it. This paper provides that, at the scale of one machine, one course, and one run per model, and with its limitations stated.
